% ============================================================================
% Chapter 2: Complete Examples and Exercises
% ============================================================================
\chapter{Complete Examples and Exercises}
\label{ch:examples-exercises}

This chapter provides hands-on examples and exercises to practice the concepts from Chapter 1. Work through these at your own pace to reinforce your understanding.

\section{Exercises for Chapter 1}
\label{sec:ch1-exercises}

Complete these exercises to practice the concepts from Chapter 1: Introduction to Embedded Rust.

\subsection{Easy Exercise: Development Environment Setup}
\label{subsec:easy-setup}

Create a shell script that automates the development environment setup for your operating system.

\textbf{Requirements:}
\begin{itemize}
    \item Check for Rust installation
    \item Install required targets (thumbv8m.main-none-eabihf)
    \item Verify probe-rs installation
    \item Provide helpful error messages
    \item Support Linux, macOS, and Windows (WSL)
\end{itemize}

\begin{minted}{bash}
#!/bin/bash
# setup-embedded-rust.sh

echo "Setting up Embedded Rust development environment..."

# Check Rust
if ! command -v rustup &> /dev/null; then
    echo "Error: Rust not found. Install from https://rustup.rs/"
    exit 1
fi

# Install target
rustup target add thumbv8m.main-none-eabihf

# Check probe-rs
if ! command -v probe-rs &> /dev/null; then
    echo "Installing probe-rs..."
    cargo install probe-rs --features cli
fi

echo "Setup complete!"
\end{minted}

\subsection{Medium Exercise: Blink an LED}
\label{subsec:medium-blink}

Modify the bare-metal example to blink an LED. Even without hardware, write the complete code.

\textbf{Requirements:}
\begin{itemize}
    \item Add LED toggle functionality
    \item Include a counter variable that increments
    \item Add debug information to panic handler
    \item Use proper GPIO configuration
\end{itemize}

\begin{minted}{rust}
#![no_std]
#![no_main]

use core::panic::PanicInfo;

// Simple counter
static mut COUNTER: u32 = 0;

#[panic_handler]
fn panic(_info: &PanicInfo) -> ! {
    // Improved panic handler
    loop {
        cortex_m::asm::bkpt();
    }
}

#[no_mangle]
pub extern "C" fn main() -> ! {
    let mut counter = 0u32;
    
    loop {
        // Toggle LED (pseudo-code)
        // gpio.toggle_led();
        
        counter += 1;
        
        // Simple delay
        for _ in 0..100_000 {
            cortex_m::asm::nop();
        }
    }
}
\end{minted}

\subsection{Hard Exercise: Type-Safe GPIO API}
\label{subsec:hard-gpio-api}

Design a type-safe GPIO API with compile-time guarantees:

\textbf{Requirements:}
\begin{itemize}
    \item Prevent reading from output pins (compile error)
    \item Prevent writing to input pins (compile error)
    \item Track pin configuration at compile time
    \item Zero-cost abstractions
    \item Support for multiple pin modes
\end{itemize}

\begin{minted}{rust}
pub struct GpioPin<MODE> {
    _mode: PhantomData<MODE>,
}

pub struct Output;
pub struct Input;
pub struct Disabled;

impl GpioPin<Output> {
    pub fn set_high(&mut self) { /* ... */ }
    pub fn set_low(&mut self) { /* ... */ }
}

impl GpioPin<Input> {
    pub fn is_high(&self) -> bool { /* ... */ }
    pub fn is_low(&self) -> bool { /* ... */ }
}

// This compiles:
let mut pin = GpioPin::<Output>::new();
pin.set_high();

// This fails to compile:
// let input_pin = GpioPin::<Input>::new();
// input_pin.set_high(); // ERROR: no method `set_high` on `Input`
\end{minted}

\section{Solutions and Discussion}
\label{sec:exercise-solutions}

\subsection{Easy Exercise Solution}
\label{subsec:solution-easy}

The shell script should detect the operating system and provide appropriate installation commands:

\begin{minted}{bash}
#!/bin/bash
# setup-embedded-rust-complete.sh

case "$(uname -s)" in
    Linux*)
        echo "Linux detected"
        # Linux-specific setup
        ;;
    Darwin*)
        echo "macOS detected"
        # macOS-specific setup
        ;;
    *)
        echo "Unsupported OS"
        exit 1
        ;;
esac
\end{minted}

\subsection{Medium Exercise Solution}
\label{subsec:solution-medium}

Key points for blinking LED:

\begin{itemize}
    \item Configure GPIO as output first
    \item Use a delay loop or timer
    \item Toggle between set_high() and set_low()
    \item Counter tracks iterations
\end{itemize}

\subsection{Hard Exercise Solution}
\label{subsec:solution-hard}

The type-state pattern uses Rust's type system to enforce pin modes at compile time:

\begin{minted}{rust}
pub struct OutputPin {
    pin_number: u8,
}

impl OutputPin {
    pub fn set(&mut self, state: bool) {
        // Unsafe: direct register access
        unsafe {
            (*GPIO::ptr()).dout |= (state as u32) << self.pin_number;
        }
    }
}

pub struct InputPin {
    pin_number: u8,
}

impl InputPin {
    pub fn read(&self) -> bool {
        unsafe {
            ((*GPIO::ptr()).din >> self.pin_number) & 1 != 0
        }
    }
}
\end{minted}

\section{Additional Examples}
\label{sec:additional-examples}

\subsection{Minimal Complete Project}
\label{subsec:minimal-complete}

Reference implementation: \url{https://github.com/bitscrafts/EFR32MG2X-RS/tree/main/examples/efr32mg24-minimal}

\subsection{Testing Framework Setup}
\label{subsec:testing-setup}

Example test structure:\\[2em]
For solutions and additional examples, visit the GitHub repository:\\[1em]
\url{https://github.com/bitscrafts/EFR32MG2X-RS}

